% Generated by experiments/run_report_artifacts.py; do not edit by hand.
% Result command: python -u experiments/run_report_artifacts.py results --only long_vertical_transition
% Result file: results/long_vertical_transition_summary.csv
\begin{table}[h]
\centering
\small
\begin{tabular}{rrrrrrrr}
\toprule
Initial tilt & Benefit & Gravity & Vertical & Collapse & Recovered & Success & Final tilt \\
\midrule
0.000 & 0.150 & 0/10 & 0/10 & 0/10 & 0/10 & 0.199 & 0.118 \\
0.000 & 0.200 & 0/10 & 0/10 & 0/10 & 0/10 & 0.193 & 0.121 \\
0.000 & 0.250 & 1/10 & 1/10 & 1/10 & 0/10 & 0.175 & 0.211 \\
0.500 & 0.150 & 3/10 & 3/10 & 3/10 & 0/10 & 0.136 & 0.366 \\
0.500 & 0.200 & 4/10 & 6/10 & 6/10 & 1/10 & 0.099 & 0.599 \\
0.500 & 0.250 & 5/10 & 9/10 & 9/10 & 1/10 & 0.037 & 0.838 \\
1.000 & 0.150 & 7/10 & 9/10 & 7/10 & 1/10 & 0.079 & 0.829 \\
1.000 & 0.200 & 5/10 & 10/10 & 10/10 & 1/10 & 0.020 & 0.899 \\
1.000 & 0.250 & 8/10 & 9/10 & 9/10 & 4/10 & 0.095 & 0.829 \\
\bottomrule
\end{tabular}
\caption{Long axial-orientation transition experiment. Gravity counts seeds in
which both mean sender and receiver transposition reached 0.50 by generation
120. Vertical counts seeds with final mean comb tilt at least 0.80. Collapse
counts seeds whose mean foraging success fell to 0.02 or below at any
generation. Recovered counts collapsed seeds whose final success was at least
0.10. Success and final tilt are final-generation means over all ten seeds.}
\label{tab:long-transition}
\end{table}
